\section*{Chapitre \ref{ch:fm:finance-product-control-and-tax} --- Finance, product control et fiscalité}

\begin{solution}{exo:fm:finance-product-control-and-tax:1}
$475\,100+63\,910+1\,580-11\,670-149\,787=379\,133$, sans rien d'inexpliqué.
\end{solution}

\begin{solution}{exo:fm:finance-product-control-and-tax:2}
$0.08/(2\times0.002)=20$ fois par an~; la totalité à 40.
\end{solution}

\begin{solution}{exo:fm:finance-product-control-and-tax:3}
$0.6\times20~\%+0.4\times37~\%=26.8~\%$.
\end{solution}

\begin{solution}{exo:fm:finance-product-control-and-tax:4}
Le total a bougé de 95\,967~\$, dans la tolérance de 100\,000~\$, car un gain sur les mouvements de marché depuis le flash compense la plus grande
partie de l'ajustement de valorisation~; la catégorie a bougé de 149\,787~\$. Les deux règles sont nécessaires~: des catégories qui se compensent masquent précisément les problèmes que le passage est censé révéler.
\end{solution}

\begin{solution}{exo:fm:finance-product-control-and-tax:5}
Le rachat intervient dans les 30 jours suivant la vente, donc la déduction de la perte est refusée~; la règle ne s'applique pas à un dealer en titres
qui réalise la perte dans le cours normal de son activité.
\end{solution}

\begin{solution}{exo:fm:finance-product-control-and-tax:6}
Ses demandes pour dol l'obligeaient à prouver qu'elle avait payé parce qu'elle avait été induite en erreur~; le tribunal a jugé que ses contrôles étaient si limités qu'elle avait payé
sans se fier aux déclarations, si bien que le lien de causalité faisait défaut, bien que les demandes de remboursement soient infondées.
\end{solution}

\begin{solution}{exo:fm:finance-product-control-and-tax:7}
Un flash de 529\,950~\$ et un définitif de 463\,393~\$~; le total bouge de 66\,557~\$, dans la tolérance, et les ajustements de valorisation dépassent à nouveau
100\,000~\$.
\end{solution}

\begin{solution}{exo:fm:finance-product-control-and-tax:8}
Un petit total peut masquer de grandes catégories qui se compensent, et un flash construit à partir des mêmes valorisations que le définitif s'accorde avec lui quelle que soit la valeur de ces valorisations.
Le test des valorisations est la vérification indépendante, non le flash.
\end{solution}

\begin{solution}{pb:fm:finance-product-control-and-tax:1}
\begin{enumerate}
\item Voir les \cref{def:fm:finance-product-control-and-tax:flash,def:fm:finance-product-control-and-tax:walk}.
\item Chaque étape est la différence de deux totaux de livres, et deux étapes consécutives partagent un total, si bien que la somme se télescope du flash au définitif.
\item Voir le \cref{tab:fm:finance-product-control-and-tax:walk}.
\item Seule la règle sur l'ajustement de valorisation~: 149\,787~\$ contre 100\,000~\$~; le total, 95\,967~\$, est dans la tolérance.
\item 73\,000 actions ont clôturé à 123.49~\$, 1.2~\% au-dessus d'un consensus de 122.01~\$, soit au bord de la bande de 0.5~\%, à 122.63~\$.
\item Une ligne hiérarchique extérieure aux desks, la propriété des valorisations officielles, et des ajustements que le trader ne peut pas annuler.
\item Que les valorisations du desk sont optimistes, que sa saisie des transactions est tardive, ou que ses frais ont été oubliés.
\item Qui a signé, quand, sur quelles preuves, chaque exception et sa raison~; les entrées ne sont qu'ajoutées.
\item Voir la \cref{def:fm:finance-product-control-and-tax:drag}.
\item Voir la \cref{prop:fm:finance-product-control-and-tax:half}.
\item 40, 20 et 8 fois par an.
\item Voir le \cref{tab:fm:finance-product-control-and-tax:tax}~: arbitrage statistique 11.85, $-8.00$ et 3.85~\%~; momentum 8.77, 6.00 et 7.58~\%~; value
6.48, 6.60 et 6.31~\%.
\item Voir l'\cref{dat:fm:finance-product-control-and-tax:tax}.
\item Voir la \cref{def:fm:finance-product-control-and-tax:cumex}.
\item 4\,170 demandes payées de mi-2012 à mi-2015, un peu moins de 12.1~milliards de DKK, aucune valide~; les demandes civiles contestées ont échoué car
l'administration n'a pas été montrée induite en erreur.
\item Obtenir l'approbation de ses fonctions fiscale et juridique, et auditer le profit comme une position fiscale.
\item Voir la \cref{def:fm:finance-product-control-and-tax:tp}~; frais de recherche et frais d'exécution.
\item Les \omterm{def:fm:finance-product-control-and-tax:tp}{prix de transfert} doivent suivre là où se trouvent le travail et les décisions~; si les documents ne concordent pas, la tarification n'a aucun fondement.
\item $T_{1/2}=r/(2t)$~; une différence au-delà de la tolérance globale, ou toute catégorie au-delà de sa propre tolérance, ou tout élément inexpliqué.
\item Valider chaque jour seulement les passages qui s'additionnent et franchissent les seuils du total et des catégories, en remontant chaque exception au registre~; n'exécuter aucune stratégie dont le
rendement dépend d'une position fiscale sans approbation fiscale et juridique préalable.
\end{enumerate}
\end{solution}

\begin{solution}{iq:fm:finance-product-control-and-tax:1}
Le flash est l'estimation précoce du desk, à partir de ses propres valorisations et des transactions saisies jusqu'à une heure limite~; le définitif inclut les transactions tardives, les frais et
les valorisations vérifiées de manière indépendante, et il est validé.
\iqlookfor{heure limite, frais et vérification.}
\end{solution}

\begin{solution}{iq:fm:finance-product-control-and-tax:2}
On construit chaque étape comme la différence de deux réévaluations du même livre, en unités exactes, et l'on signale tout reliquat comme inexpliqué.
\iqlookfor{réévaluations qui se télescopent.}
\end{solution}

\begin{solution}{iq:fm:finance-product-control-and-tax:3}
L'impôt coûte $25\times0.20~\%=5~\%$ par an, il reste 1~\%~: les cinq sixièmes du rendement vont à l'impôt. Pas dans cet instrument~; on négocie là où
l'impôt ne s'applique pas, si les règles le permettent, ou l'on ralentit la stratégie.
\iqlookfor{rotation multipliée par le taux, comparée au rendement.}
\end{solution}

\begin{solution}{iq:fm:finance-product-control-and-tax:4}
On le suit comme un schéma, on resserre les bandes ou l'on exige des prix observables, on applique des réserves, et on le porte auprès du responsable du desk et du risque~:
les valorisations sont l'estimation, par le trader, de sa propre rémunération.
\iqlookfor{indépendance et schémas sur plusieurs jours.}
\end{solution}

\begin{solution}{iq:fm:finance-product-control-and-tax:5}
La section 1256 impose les gains sur contrats à terme à 60~\% en long terme et 40~\% en court terme, quelle que soit la durée de détention~; les mêmes gains sur des actions détenues pendant des semaines sont tous
à court terme.
\iqlookfor{la règle 60/40.}
\end{solution}

\begin{solution}{iq:fm:finance-product-control-and-tax:6}
Seuils sur le total et sur chaque catégorie, ajustés à la volatilité et à la taille de chaque desk, une tolérance nulle sur le P\&L inexpliqué, des alertes de tendance sur
les catégories, et des voies de remontée avec des délais.
\iqlookfor{ajustés, par catégorie et suivis dans le temps.}
\end{solution}
